\section{De los modelos de lenguaje a los asistentes}

\subsection{Limitaciones de los modelos de lenguaje preentrenados}

Los modelos de lenguaje preentrenados aprenden a \textbf{completar texto}, pero eso no equivale a \textbf{ayudar al usuario}. Entre el objetivo de modelado del lenguaje y «satisfacer las preferencias humanas» existe un desalineamiento de fondo.

\begin{knowledgebox}{Qué se aprende en el preentrenamiento}
Con grandes cantidades de texto, el preentrenamiento adquiere un conocimiento amplio: hechos, gramática, resolución de correferencias, análisis de sentimiento, razonamiento básico, matemáticas, programación, etc. Sin embargo, lo que aprende es «el siguiente token más probable», no «la respuesta que el usuario más desea».
\end{knowledgebox}

\subsection{La diversidad de las preferencias humanas}

Las preferencias reales de los usuarios no forman una sola función de recompensa: se componen de varias dimensiones, como el grado de cumplimiento de la tarea, la creatividad, la cortesía, la interpretabilidad e incluso la estética. En la clase, Archit afirma que \textit{«las preferencias son un espacio vectorial multimodal, no un escalar único»}, y subraya que, al construir las señales, es indispensable recopilarlas por niveles.

Según su origen, las preferencias pueden clasificarse así:
\begin{itemize}
    \item \textbf{Comparaciones humanas}: el usuario o el anotador recibe una lista de respuestas y elige la mejor.
    \item \textbf{Métricas verificables}: los problemas de matemáticas, las tareas de programación, la coherencia factual, etc., pueden calificarse de forma automática.
    \item \textbf{Datos de comportamiento}: la tasa de clics, la tasa de reescritura, la tasa de abandono, etc., reflejan de manera implícita la satisfacción del usuario.
    \item \textbf{Restricciones de gobernanza}: los requisitos de seguridad y de cumplimiento normativo exigen señales adicionales de guardrail.
\end{itemize}

\begin{knowledgebox}{Idea central del modelado con matriz de preferencias}
Combinar varias señales de preferencia en una \textit{hierarchical preference matrix} permite delimitar con claridad el alcance de cada señal y alinear objetivos distintos mediante una ponderación durante la optimización.
\end{knowledgebox}

\begin{importantbox}{Interpretabilidad de las preferencias con múltiples señales}
Dividir las signal en dos clases, \textit{primary preference} (satisfacción central) y \textit{secondary constraint} (seguridad, bias), y tratarlas por separado con un reward model y con un guardrail filter, respectivamente, reduce el reward hacking y aumenta la transparencia de la gobernanza.
\end{importantbox}

\subsection{Resumen de la sección}

Un LM preentrenado posee un conocimiento amplio, pero para convertirse en asistente debe alinearse en varias dimensiones, como el fine-tuning con instrucciones, la separación de las señales de preferencia y la incorporación de señales de gobernanza; solo así se evita que una reward única lo desvíe.
